\chapter{Context of the Project}
\chaptertoc

\section{Introduction}
In this chapter, we present the general context of the project. We begin with the host organization and its sectors of activity. Then, we introduce the project presentation, including the problem statement, a brief study of existing approaches, and the proposed solution. Finally, we describe the methodology adopted to guide the work throughout the project.

\section{Host Organization}

\subsection{DataDoIt}
This final-year project was carried out in collaboration with DataDoIt, which provided the professional framework, guidance, and practical use-case context.

\subsection{Sectors of Activity}
The problem addressed by this work is relevant to several sectors where queue congestion impacts service quality, including retail checkout systems, banking agencies, and customer service centers.

\section{Presentation of the Project}

\subsection{Project Context}
This project focuses on real-time queue monitoring from video streams in order to transform raw visual data into operational indicators for decision support.

\subsection{Problem Statement}
Manual approaches or simple counting-based methods remain limited:
\begin{itemize}
    \item they are time-consuming,
    \item they do not explicitly estimate waiting time,
    \item they do not report uncertainty in the measurements.
\end{itemize}

\subsection{Study of Existing Approaches}

\subsubsection{Study of Existing and Similar Solutions}
Existing solutions are generally based on manual counting, basic counters, or isolated analytics components that do not provide a full end-to-end queue estimation pipeline.

\subsubsection{Limitations of Existing Solutions}
Most existing approaches either ignore uncertainty, provide delayed analytics, or fail to combine detection, tracking, queue modeling, and actionable alerting in one integrated system.

\subsection{Proposed Solution}
The proposed system combines person detection, multi-object tracking, zone-based filtering, queue-theory metrics, and Bayesian uncertainty quantification with dashboard and webhook integration.

\subsection{Objectives of the Project}
This work aims to design a system able to:
\begin{enumerate}
    \item detect and track people in real time,
    \item count customers inside a configurable queue zone,
    \item estimate arrival rate $\lambda$ and service rate $\mu$,
    \item compute expected waiting time $W$,
    \item quantify uncertainty for key metrics,
    \item provide a monitoring interface and external integration (n8n).
\end{enumerate}

\subsection{Target Audience}
The main audience includes operational staff (for real-time supervision), managers (for resource planning), and end-users who benefit from reduced waiting time and improved service flow.

\section{Project Management Methodology}

\subsection{Choice of Methodology: Agile}
The project follows an incremental and iterative approach to deliver functional modules progressively while integrating feedback and resolving technical risks early.

\subsection{SCRUM}
Development was organized in explicit Scrum sprints with incremental deliverables: Sprint 1 (core perception pipeline), Sprint 2 (GUI and multi-video support), Sprint 3 (YOLO26 migration and performance tuning), Sprint 4 (uncertainty quantification), and Sprint 5 (n8n integration and alerting).

\section{Conclusion}
This chapter introduced the project context, organizational setting, problem statement, and proposed approach. The next chapter presents the state of the art supporting the technical choices made in this work.
